\chapter{Verwandte Arbeiten}
\label{chapter:related-work}

\section{Grafana als Visualisierungsschicht für Monitoring-Daten}
\label{sec:related-grafana}

Der Einsatz von \gls{grafana} als zentrale Visualisierungsschicht für heterogene
Monitoring-Daten wurde bereits in unterschiedlichen Anwendungskontexten untersucht.
Bondar et al. beschreiben den Einsatz von \gls{grafana} als Präsentationsschicht für das
Kontrollsystem Karabo am European XFEL, wo \gls{grafana}-Dashboards eine nahezu
echtzeitfähige Überwachung von Steuerungsdaten ermöglichen und dabei verschiedene
Monitoring-Szenarien sowie gewonnene Erkenntnisse aus dem praktischen Einsatz
diskutiert werden~\cite{VBondar2025}. Diese Arbeit bestätigt die Eignung von
\gls{grafana} für zeitkritische Monitoring-Anwendungen, wie sie auch der in dieser
Arbeit beschriebenen Reaktionszeitanforderung zugrunde liegt (siehe
Abschnitt~\ref{subsec:reaction-times}), allerdings in einem wissenschaftlichen statt
industriellen Kontext und ohne die Notwendigkeit, mehrere unabhängige, proprietäre
Drittanbieter-Systeme über ein eigenständiges Plugin anzubinden.

\section{Integration heterogener Datenquellen}
\label{sec:related-integration}

Die Zusammenführung von Daten aus unterschiedlichen, technisch heterogenen Systemen
stellt eine zentrale Herausforderung dar, die auch außerhalb des Monitoring-Kontexts
untersucht wird. Putrama und Martinek geben einen umfassenden Überblick über
bestehende Ansätze zur Integration heterogener Datenquellen und unterscheiden dabei
insbesondere zwischen physischer und virtueller Datenintegration. Sie stellen fest,
dass physische Integrationssysteme zwar über eine leistungsfähigere
Abfrageverarbeitung verfügen, jedoch mit höherem Implementierungs- und
Wartungsaufwand verbunden sind, während virtuelle Integrationsansätze zunehmend an
Bedeutung gewinnen~\cite{Putrama2024}. Diese Erkenntnis deckt sich mit der in
Abschnitt~\ref{subsec:reaction-times} getroffenen Entscheidung für eine virtuelle
Datenintegration innerhalb des entwickelten Plugins, bei der die vorliegende Arbeit
den beschriebenen Zusammenhang zwischen Implementierungsaufwand und Antwortzeit im
konkreten Anwendungsfall mehrerer Netzwerk- und Systemüberwachungsquellen praktisch
bestätigt.

\section{Zentrale Darstellung verteilter Infrastruktur}
\label{sec:related-spog}

Das Konzept einer zentralen, konsolidierten Sicht auf verteilte, heterogene
Infrastruktur wird in der Literatur unter anderem im Kontext des sogenannten
\enquote{Single Pane of Glass}-Ansatzes diskutiert. Velez-Rojas et al. weisen darauf
hin, dass Netzwerk-Verantwortliche in modernen Unternehmensumgebungen mit sehr
großen Mengen an Informationen umgehen müssen, und schlagen eine Visualisierungsmethode
vor, die Übersichten generiert und dabei die für die jeweilige Nutzerin oder den
jeweiligen Nutzer relevantesten Zusammenhänge hervorhebt~\cite{VelezRojas2013}. Der
in dieser Arbeit entwickelte Ansatz verfolgt ein ähnliches Ziel auf Ebene einzelner
Standorte: Anstatt sämtliche Detailinformationen gleichwertig darzustellen, werden,
wie in Abschnitt~\ref{sec:systemmetriken} und Abschnitt~\ref{sec:netzwerkmetriken}
beschrieben, zunächst aggregierte Standortwerte präsentiert, die bei Bedarf eine
gezielte Detailbetrachtung einzelner Systeme ermöglichen.